\chapter{The Axioms of Probability}
\label{ch:axioms}

The classical rule $\abs{A}/\abs{\Omega}$ works only for finite, symmetric
situations. In 1933 Kolmogorov showed that all of probability can be built on a
short list of axioms. This chapter states them, derives their consequences, and
proves the two most useful computational tools: inclusion--exclusion and
continuity of probability.

\section{Events as a $\sigma$-algebra}

When $\Omega$ is finite or countable, every subset can be an event. In general
we need to choose a collection $\F$ of subsets of $\Omega$ that are allowed to
be events, and this collection must be closed under the operations we use to
combine events.

\begin{definition}[$\sigma$-algebra]
A collection $\F$ of subsets of $\Omega$ is a \emph{$\sigma$-algebra} (or
$\sigma$-field) if
\begin{enumerate}[label=(\roman*)]
\item $\Omega\in\F$;
\item if $A\in\F$ then $A^c\in\F$;
\item if $A_1,A_2,\dots\in\F$ then $\bigcup_{i=1}^\infty A_i\in\F$.
\end{enumerate}
\end{definition}

By De Morgan's laws (Theorem~\ref{thm:demorgan}), $\F$ is then also closed under
countable intersections, and $\emptyset=\Omega^c\in\F$. Taking
$A_{n+1}=A_{n+2}=\dots=\emptyset$ shows that $\F$ is closed under finite unions
as well.

\begin{nuance}{why not simply use all subsets?}
For finite or countable $\Omega$ we can, and we will, take $\F=2^\Omega$, the
collection of all subsets. For $\Omega=[0,1]$ one cannot assign a length
consistently to \emph{every} subset: there exist sets (Vitali sets) for which
no assignment of ``length'' is compatible with countable additivity and
translation invariance. The $\sigma$-algebra restricts attention to the sets
that can be measured. This is a technicality that rarely matters in computation;
we return to it in Chapter~\ref{ch:measure}. For now, read ``event'' as ``measurable subset''
and note that every set built from intervals by countable unions,
intersections and complements is an event.
\end{nuance}

\section{Kolmogorov's axioms}

\begin{definition}[Probability measure]
Let $\Omega$ be a set and $\F$ a $\sigma$-algebra on $\Omega$. A
\emph{probability measure} is a function $\Prob:\F\to\R$ such that
\begin{enumerate}[label=\textbf{A\arabic*.}]
\item (\emph{Non-negativity}) $\Prob(A)\ge 0$ for every $A\in\F$;
\item (\emph{Normalisation}) $\Prob(\Omega)=1$;
\item (\emph{Countable additivity}) if $A_1,A_2,\dots\in\F$ are pairwise
disjoint, then
\[
\Prob\Bigl(\bigcup_{i=1}^\infty A_i\Bigr)=\sum_{i=1}^\infty\Prob(A_i).
\]
\end{enumerate}
The triple $(\Omega,\F,\Prob)$ is called a \emph{probability space}.
\end{definition}

These axioms say nothing about how to choose $\Prob$ in a given problem; that is
a modelling decision. They only say what any sensible assignment must satisfy.

\begin{example}[Discrete spaces]\label{ex:discrete}
Let $\Omega=\{\omega_1,\omega_2,\dots\}$ be finite or countable and let
$p_1,p_2,\dots\ge 0$ satisfy $\sum_ip_i=1$. Defining
$\Prob(A)=\sum_{i:\omega_i\in A}p_i$ gives a probability measure on
$\F=2^\Omega$: non-negativity and normalisation are immediate, and countable
additivity follows because a sum of non-negative terms may be rearranged and
grouped freely. The classical definition is the case $p_i=1/\abs\Omega$.
\end{example}

\section{Elementary consequences}

\begin{proposition}\label{prop:basic}
Let $(\Omega,\F,\Prob)$ be a probability space and $A,B\in\F$. Then
\begin{enumerate}[label=(\roman*)]
\item $\Prob(\emptyset)=0$;
\item (finite additivity) if $A\cap B=\emptyset$ then
$\Prob(A\cup B)=\Prob(A)+\Prob(B)$; more generally this holds for any finite
number of pairwise disjoint events;
\item $\Prob(A^c)=1-\Prob(A)$;
\item (monotonicity) if $A\subseteq B$ then $\Prob(A)\le\Prob(B)$, and
$\Prob(B\setminus A)=\Prob(B)-\Prob(A)$;
\item $0\le\Prob(A)\le 1$;
\item $\Prob(A\cup B)=\Prob(A)+\Prob(B)-\Prob(A\cap B)$.
\end{enumerate}
\end{proposition}
\begin{proof}
(i) Apply A3 to the sequence $\emptyset,\emptyset,\dots$ (pairwise disjoint,
with union $\emptyset$): $\Prob(\emptyset)=\sum_{i=1}^\infty\Prob(\emptyset)$.
A real number $x$ with $x=\sum_{i\ge 1}x$ must be $0$, because $x>0$ would give
$+\infty$ and $x\ge0$ by A1.

(ii) Apply A3 to $A,B,\emptyset,\emptyset,\dots$ and use (i). The case of
$m$ events is identical.

(iii) $A$ and $A^c$ are disjoint with union $\Omega$, so by (ii) and A2,
$1=\Prob(A)+\Prob(A^c)$.

(iv) If $A\subseteq B$, then $B=A\cup(B\setminus A)$ is a disjoint union, so
$\Prob(B)=\Prob(A)+\Prob(B\setminus A)\ge\Prob(A)$ by A1.

(v) The upper bound is (iv) with $B=\Omega$; the lower bound is A1.

(vi) Write $A\cup B=A\cup(B\setminus A)$, a disjoint union, and
$B=(A\cap B)\cup(B\setminus A)$, also disjoint. Then
$\Prob(A\cup B)=\Prob(A)+\Prob(B\setminus A)$ and
$\Prob(B)=\Prob(A\cap B)+\Prob(B\setminus A)$. Subtract the second equation
from the first.
\end{proof}

\begin{theorem}[Boole's inequality / union bound]\label{thm:boole}
For any events $A_1,A_2,\dots$ (finitely or countably many, not necessarily
disjoint),
\[
\Prob\Bigl(\bigcup_iA_i\Bigr)\le\sum_i\Prob(A_i).
\]
\end{theorem}
\begin{proof}
Define $B_1=A_1$ and $B_i=A_i\setminus(A_1\cup\dots\cup A_{i-1})$ for
$i\ge 2$. The $B_i$ are pairwise disjoint, $B_i\subseteq A_i$, and
$\bigcup_iB_i=\bigcup_iA_i$. (If $\omega\in\bigcup_iA_i$, let $i$ be the
smallest index with $\omega\in A_i$; then $\omega\in B_i$.) By countable
additivity and monotonicity,
\[
\Prob\Bigl(\bigcup_iA_i\Bigr)=\sum_i\Prob(B_i)\le\sum_i\Prob(A_i).\qedhere
\]
\end{proof}

The union bound requires no independence assumption, which makes it one of the
most widely applicable inequalities in probability.

\section{Inclusion--exclusion}

Part (vi) of Proposition~\ref{prop:basic} computes $\Prob(A\cup B)$ for two
events. For $n$ events the formula alternates between adding and subtracting
intersections.

\begin{theorem}[Inclusion--exclusion]\label{thm:incexc}
For events $A_1,\dots,A_n$ in a discrete probability space,
\[
\Prob\Bigl(\bigcup_{i=1}^nA_i\Bigr)=\sum_{k=1}^n(-1)^{k+1}S_k,
\qquad
S_k=\sum_{1\le i_1<\dots<i_k\le n}\Prob\bigl(A_{i_1}\cap\dots\cap A_{i_k}\bigr).
\]
\end{theorem}
\begin{proof}
We give the proof for the discrete setting of Example~\ref{ex:discrete}, where
$\Prob(A)=\sum_{\omega\in A}p(\omega)$ (the result holds on any probability
space, and a proof by induction on $n$ using Proposition~\ref{prop:basic}(vi)
is possible as well). It suffices to show that every outcome $\omega$
contributes the same weight $p(\omega)$ to both sides, or contributes nothing
to both.

If $\omega$ lies in none of the $A_i$, it contributes nothing to either side.
Otherwise suppose $\omega$ lies in exactly $m\ge1$ of the events. On the left
it contributes $p(\omega)$. On the right, $\omega$ belongs to the intersection
$A_{i_1}\cap\dots\cap A_{i_k}$ exactly when all of $i_1,\dots,i_k$ are among the
$m$ events containing $\omega$; there are $\binom mk$ such index sets. So the
total contribution of $\omega$ to the right side is
\[
p(\omega)\sum_{k=1}^m(-1)^{k+1}\binom mk
=p(\omega)\Bigl[1-\sum_{k=0}^m(-1)^k\binom mk\Bigr]=p(\omega)(1-0)=p(\omega),
\]
by Corollary~\ref{cor:sums}. The two sides agree.
\end{proof}

\begin{theorem}[Bonferroni inequalities]\label{thm:bonferroni}
With $S_k$ as above, the partial sums alternately overestimate and
underestimate:
\[
\sum_{k=1}^{j}(-1)^{k+1}S_k\ \ge\ \Prob\Bigl(\bigcup_iA_i\Bigr)\ \text{ for odd } j,
\qquad
\sum_{k=1}^{j}(-1)^{k+1}S_k\ \le\ \Prob\Bigl(\bigcup_iA_i\Bigr)\ \text{ for even } j.
\]
\end{theorem}
\begin{proof}
We use the same outcome-by-outcome argument. First an identity: for
integers $m\ge1$ and $0\le j\le m-1$,
\begin{equation}\label{eq:partial}
\sum_{k=0}^{j}(-1)^k\binom mk=(-1)^j\binom{m-1}{j}.
\end{equation}
Proof by Pascal's rule (Theorem~\ref{thm:pascal}), writing
$\binom mk=\binom{m-1}k+\binom{m-1}{k-1}$ for $k\ge 1$: the sum becomes
$\sum_{k=0}^j(-1)^k\binom{m-1}k+\sum_{k=1}^j(-1)^k\binom{m-1}{k-1}$, and the
second sum equals $-\sum_{k=0}^{j-1}(-1)^k\binom{m-1}{k}$. The two sums cancel
except for the term $k=j$ of the first, which is $(-1)^j\binom{m-1}j$.

Now consider an outcome lying in exactly $m\ge 1$ of the events. Its
contribution to the partial sum up to $j$ is, by \eqref{eq:partial} (and
trivially if $j\ge m$, where the full alternating sum applies),
\[
p(\omega)\sum_{k=1}^{j}(-1)^{k+1}\binom mk=p(\omega)\Bigl[1-(-1)^j\binom{m-1}{j}\Bigr].
\]
For $j$ odd the bracket equals $1+\binom{m-1}{j}\ge 1$, and for $j$ even it
equals $1-\binom{m-1}{j}\le1$. Summing over $\omega$ gives the claims.
\end{proof}

\begin{example}[The matching problem]\label{ex:matching}
$n$ letters are put at random into $n$ addressed envelopes, one per envelope.
What is the probability that no letter lands in its own envelope? Let
$\Omega$ be the set of $n!$ permutations, all equally likely, and let $A_i$ be
the event that letter $i$ is correctly placed (position $i$ is fixed). For any
$k$ distinct indices, the number of permutations fixing all of them is
$(n-k)!$, so
\[
\Prob(A_{i_1}\cap\dots\cap A_{i_k})=\frac{(n-k)!}{n!},
\qquad
S_k=\binom nk\frac{(n-k)!}{n!}=\frac1{k!}.
\]
By inclusion--exclusion,
\[
\Prob(\text{no match})=1-\Prob\Bigl(\bigcup_iA_i\Bigr)
=1-\sum_{k=1}^n\frac{(-1)^{k+1}}{k!}=\sum_{k=0}^n\frac{(-1)^k}{k!}.
\]
This is the partial sum of the series $e^{-1}=\sum_{k\ge 0}(-1)^k/k!$. By the
alternating series estimate the error is at most $1/(n+1)!$, so already for
$n=7$ the answer agrees with $e^{-1}\approx0.3679$ to four decimal places.
Remarkably, the limit does not depend on $n$.
\end{example}

\begin{interview}{inclusion--exclusion versus complement}
When asked for the probability that \emph{none} of several events occurs, and
the events are symmetric (each $k$-fold intersection has probability depending
only on $k$), inclusion--exclusion gives a closed form quickly, as in the
matching problem. The same structure appears in the coupon-collector and
``everyone has a different birthday'' problems.
\end{interview}

\section{Continuity of probability}

Countable additivity is the axiom that connects probability to limits. Its most
useful form is continuity.

\begin{definition}
A sequence of events $A_1,A_2,\dots$ is \emph{increasing} if
$A_1\subseteq A_2\subseteq\cdots$, with limit $A=\bigcup_nA_n$ (written
$A_n\uparrow A$). It is \emph{decreasing} if $A_1\supseteq A_2\supseteq\cdots$,
with limit $A=\bigcap_nA_n$ (written $A_n\downarrow A$).
\end{definition}

\begin{theorem}[Continuity of probability]\label{thm:continuity}
If $A_n\uparrow A$ or $A_n\downarrow A$, then $\Prob(A_n)\to\Prob(A)$.
\end{theorem}
\begin{proof}
\emph{Increasing case.} Set $B_1=A_1$ and $B_n=A_n\setminus A_{n-1}$ for
$n\ge 2$. The $B_n$ are pairwise disjoint, $A_n=B_1\cup\dots\cup B_n$, and
$A=\bigcup_nB_n$. By countable additivity,
\[
\Prob(A)=\sum_{n=1}^\infty\Prob(B_n)=\lim_{N\to\infty}\sum_{n=1}^N\Prob(B_n)
=\lim_{N\to\infty}\Prob(A_N),
\]
where the last step is finite additivity.

\emph{Decreasing case.} If $A_n\downarrow A$ then $A_n^c\uparrow A^c$ (by De
Morgan's laws). By the increasing case, $\Prob(A_n^c)\to\Prob(A^c)$, that is,
$1-\Prob(A_n)\to1-\Prob(A)$.
\end{proof}

\begin{nuance}{the decreasing case needs a finite total mass}
The decreasing case used $\Prob(\Omega)=1<\infty$. For measures with infinite
total mass (for example length on $\R$) continuity from above can fail:
$A_n=[n,\infty)$ decreases to $\emptyset$ but each has infinite length.
Probability measures escape this problem precisely because of normalisation.
\end{nuance}

\begin{example}[An infinite run of heads]\label{ex:allheads}
Toss a fair coin repeatedly. Let $A_n$ be the event that the first $n$ tosses
are all heads. Then $\Prob(A_n)=2^{-n}$ and $A_n\downarrow A$, where $A$ is the
event that \emph{every} toss is heads. By continuity,
$\Prob(A)=\lim_n2^{-n}=0$. The sequence $HHH\cdots$ is a perfectly possible
outcome, yet it has probability zero.
\end{example}

\begin{nuance}{probability zero does not mean impossible}
Example~\ref{ex:allheads} shows that an event can be possible and still have
probability $0$. Likewise, probability $1$ does not mean certain. We say an
event holds \emph{almost surely} (a.s.) if it has probability $1$. This
distinction becomes unavoidable for continuous random variables, where every
single value has probability zero but some value must occur.
\end{nuance}

\begin{nuance}{why countable, and not merely finite, additivity}
Finite additivity gives no control over limits of events; countable additivity
is exactly what makes Theorem~\ref{thm:continuity} true. It also rules out
some tempting models. There is no way to choose a natural number ``uniformly at
random'': if $\Prob(\{n\})=c$ for all $n\in\N$, then countable additivity gives
$1=\Prob(\N)=\sum_{n}c$, which equals $0$ if $c=0$ and $\infty$ if $c>0$,
a contradiction either way. (If only finite additivity is demanded, uniform
``distributions'' on $\N$ do exist, but they fail to behave well under limits,
which is why the standard theory does not allow them.)
\end{nuance}

\begin{corollary}[Countable unions as limits of finite unions]
For any events $A_1,A_2,\dots$ we have
$\Prob(\bigcup_{i=1}^\infty A_i)=\lim_{N\to\infty}\Prob(\bigcup_{i=1}^NA_i)$.
\end{corollary}
\begin{proof}
The sets $C_N=\bigcup_{i=1}^NA_i$ increase to $\bigcup_iA_i$; apply
Theorem~\ref{thm:continuity}.
\end{proof}

\section*{Exercises}
\addcontentsline{toc}{section}{Exercises}

\begin{enumerate}
\item Show that if $\Prob(A)=0.6$ and $\Prob(B)=0.7$, then
$0.3\le\Prob(A\cap B)\le0.6$, and show that both bounds can be attained.
\item Prove the Bonferroni-type bound
$\Prob(\bigcap_iA_i)\ge1-\sum_i\Prob(A_i^c)$.
\item Prove inclusion--exclusion for three events by repeated use of
Proposition~\ref{prop:basic}(vi), and then prove the general theorem by induction
on $n$.
\item (Derangements) Let $D_n$ be the number of permutations of $\{1,\dots,n\}$
with no fixed point. Show $D_n=n!\sum_{k=0}^n(-1)^k/k!$ and verify that
$D_n$ is the nearest integer to $n!/e$ for $n\ge 1$.
\item Show that the probability that exactly one of the events $A_1,\dots,A_n$
occurs equals $S_1-2S_2+3S_3-\cdots+(-1)^{n+1}nS_n$.
\item Let $\F$ be a $\sigma$-algebra. Prove that $\F$ is closed under countable
intersections and under set difference.
\item A fair coin is tossed repeatedly. Use continuity of probability to show
that the probability that heads never appears is $0$, and that the probability
that the first head appears at some finite toss is $1$.
\item Show that a function satisfying A1, A2 and finite additivity (but not
countable additivity) is a probability measure if and only if it is continuous
at $\emptyset$, meaning $\Prob(A_n)\to0$ whenever $A_n\downarrow\emptyset$.
\item (Challenge) Show that there is no probability measure on $(\N,2^\N)$ with
$\Prob(\{n\})=c$ for all $n$, but that for $p_n=2^{-n}$ the construction of
Example~\ref{ex:discrete} gives a legitimate geometric model for the first
head in repeated tosses (Example~\ref{ex:spaces}(d)).
\end{enumerate}
